% =====================================================================
% todo: check the report section
\newpage
\section{Результаты работы программы}
\label{sec:results}

Ниже приведены примеры работы программы (вывод в консоль вместе с введенными пользователем значениями). Поле выводится построчно: символ \texttt{+} обозначает живую клетку, пробел --- мертвую.

После запуска программа выводит заголовок, запрашивает ширину поля и количество итераций, а затем предлагает выбрать способ задания начальных условий (рис.~\ref{fig:start}). В квадратных скобках указаны значения по умолчанию, которые принимаются при пустом вводе.

\begin{figure}[H]
	\centering
	\console[\consolemedium]{start}
	\caption{Запуск программы и ввод параметров}
	\label{fig:start}
\end{figure}

Все значения вводятся с проверкой: если введено не целое число (\texttt{abc}, \texttt{2.5}) или число вне допустимого диапазона (\texttt{-5}, \texttt{1000}), программа выводит сообщение с допустимым диапазоном и повторяет запрос. Для \texttt{seed} допускается любое целое число или пустой ввод. На рис.~\ref{fig:error} показан пример запуска, в котором для каждого параметра сначала было введено некорректное значение, а затем корректное. После ввода всех параметров программа продолжает работу в обычном режиме.

\begin{figure}[H]
	\centering
	\console[\consolemedium]{error}
	\caption{Обработка некорректного ввода (фрагмент вывода программы)}
	\label{fig:error}
\end{figure}

При выборе случайных начальных условий пользователь задает плотность живых клеток в процентах и \texttt{seed} генератора псевдослучайных чисел; при пустом вводе \texttt{seed} генерируется автоматически и выводится на экран, чтобы запуск можно было повторить. Затем задаются частота вывода поколений и необходимость остановки при обнаружении цикла, после чего выводится начальное поколение. На рис.~\ref{fig:randinit} показано продолжение запуска с рис.~\ref{fig:start}: поле $20 \times 20$, плотность 30\,\%, $\text{seed} = 1$; в начальном поколении 118 живых клеток (ожидаемое значение $0{,}3 \cdot 400 = 120$).

\begin{figure}[H]
	\centering
	\console[\consolemedium]{randinit}
	\caption{Случайное заполнение поля (продолжение рис.~\ref{fig:start})}
	\label{fig:randinit}
\end{figure}

При ручном вводе пользователь вводит строки поля сверху вниз: символы \texttt{+} и \texttt{1} обозначают живую клетку, остальные --- мертвую; пустая строка завершает ввод, недостающие строки и столбцы считаются мертвыми (рис.~\ref{fig:manual}). В примере на поле $9 \times 9$ введен паттерн «крест».

\begin{figure}[H]
	\centering
	\console[\consolemedium]{manual}
	\caption{Ручной ввод начальной конфигурации}
	\label{fig:manual}
\end{figure}

На рис.~\ref{fig:run} показано продолжение этого запуска: вывод поколений 1--3 и итог работы. Видно, что паттерн смещается вправо, как описано в разделе~\ref{sec:analysis} (рис.~\ref{fig:cross}), и остается в строках, в которых был задан, что соответствует свойству 5 из раздела~\ref{sec:analysis}. За 3 итерации цикл не найден, поэтому итог содержит соответствующее сообщение.

\begin{figure}[H]
	\centering
	\console[\consolemedium]{run}
	\caption{Вывод поколений и итог работы (продолжение рис.~\ref{fig:manual})}
	\label{fig:run}
\end{figure}

При обнаружении повторения конфигурации программа выводит поколение, на котором найден цикл, с предпериодом и периодом в заголовке, а затем итог работы (рис.~\ref{fig:cycle}). В примере одиночная клетка в центре поля $7 \times 7$ за 7 шагов обходит тор и возвращается в исходное положение: предпериод равен 0, период --- 7. Чтобы не выводить промежуточные поколения, задано $k = 100$; поколение, на котором найден цикл, выводится всегда.

\begin{figure}[H]
	\centering
	\console[\consolemedium]{cycle}
	\caption{Обнаружение цикла и итог работы}
	\label{fig:cycle}
\end{figure}

Если цикл за заданное число итераций не найден, программа сообщает об этом (рис.~\ref{fig:nocycle}). Запуск выполнен на поле $20 \times 20$ со случайными начальными условиями (плотность 30\,\%, $\text{seed} = 1$), 100\,000 итераций и $k = 100\,000$, поэтому выведены только начальное и последнее поколения.

\begin{figure}[H]
	\centering
	\console[\consolemedium]{nocycle}
	\caption{Итог работы, цикл не обнаружен (фрагмент: последнее поколение и итог)}
	\label{fig:nocycle}
\end{figure}
